% Moved from Paper I Section 8 (governance/00b, migration map); mathematical content
% unchanged apart from cross-references, framing sentences, and the closing
% subsection on tearing.
The arithmetic motions of Section~\ref{sec:p0-aes} do not commute: addition and
multiplication transport each other's rulers by the Baumslag--Solitar-type
relation \eqref{eq:p0-bs-relation}, and the affine cocycle of
Section~\ref{sec:p0-affine-cocycles} records the resulting order dependence in
two translation coordinates.  This section constructs a commutativization of
that non-commutative motion, and the level at which commutation is imposed must
be stated exactly.

Take the history language itself: the free monoid generated by the two families
of one-step histories, additive steps and multiplicative steps.  Impose that
generators of the two families commute.  Equivalently --- and this is the form
used below --- keep the total additive charge and the total logarithmic
multiplicative charge of a history and forget the order in which they were
accumulated; the charges satisfy the imposed commutation by construction.  The
resulting object is a plane with two coordinates, and we call it the
\emph{accumulative commutative space}, ACS.  The construction has two distinct
uses, and they must not be conflated.  For a positive add--scale history it
records the two total charges; and it supplies a weighted integral along the
charge path whose value restores the
translation part that the charge endpoint alone forgets.  The ACS is therefore a
*shadow* of
the presented history: it is a commutative quotient object attached to a
non-commutative one, not another presentation of the full arithmetic syntax.

This is *not* the abelianization of the group of affine motions, and the
difference is worth recording once.  In $\operatorname{Aff}^{+}(1,\mathbb R)$
one has $D_kT_pD_k^{-1}T_p^{-1}=T_{(k-1)p}$, so the commutator subgroup is the
translation subgroup and the abelianization keeps the scale alone:
\[
  \operatorname{Aff}^{+}(1,\mathbb R)_{\mathrm{ab}}
  \;\cong\;\mathbb R_{>0}\;\cong\;(\mathbb R,+).
\]
The ACS keeps the additive charge as well, because it is attached to the
presented histories *before* evaluation rather than to their images in the
affine group.  The hierarchy --- literal histories, then charges, then evaluated
motions --- is accordingly not a convenience of exposition; collapsing it
changes the object.  Subsection~\ref{subsec:p0-acs-forgets} states what the ACS
forgets, and the distributivity example of
Subsection~\ref{subsec:p0-distributivity} already shows why the charge map
cannot be pushed through the operator quotient.

The geometric use is the reason for the name of this chapter.  Two histories
with the same charge endpoint have charge paths that
together bound a region of the ACS, and the exact value of their order defect is
a weighted area of that region (Theorem~\ref{thm:torsion-stokes}).  In
Section~\ref{subsec:p0-tearing} we read that identity as a statement about the
original AES: the arithmetic space carries a defect that cannot be seen in any
single motion, and the geometry in which it becomes an area is the
commutativized charge geometry of the presented histories.

Throughout this section an elementary additive step and an elementary scaling
step are denoted by
\[
    \mathsf A_p(x)=x+p,
    \qquad
    \mathsf M_q(x)=e^q x,
    \qquad p,q\in\mathbb R.
\]
For a history \(\gamma=(g_1,\ldots,g_n)\), the entries are performed from left
to right, so that
\[
    \nu_x(\gamma)=g_n\circ\cdots\circ g_1(x).
\]
Negative values of \(p\) and \(q\) are allowed.  Thus the same notation covers
inverse translations and inverse positive scalings.
In this section an \emph{add--scale history} means a finite word in these two
families of generators.  Such a word evaluates in
\(\operatorname{Aff}^+(1,\mathbb R)\).  The ACS construction is attached to
this presented word; it is not asserted for negative multipliers or for an
arbitrary affine syntax before that syntax has been written in add--scale form.

\subsection{The charge path}

\begin{definition}[Accumulative commutative space and charge path]
\label{def:acs}
The \emph{accumulative commutative space}, abbreviated ACS, is the oriented
plane
\[
    \mathcal A=\mathbb R^2_{(A,M)},
    \qquad \text{with orientation } dA\wedge dM.
\]
An additive step \(\mathsf A_p\) has charge increment \((p,0)\), and a scaling
step \(\mathsf M_q\) has charge increment \((0,q)\).  Starting at
\((A_0,M_0)=(0,0)\), accumulate these increments in chronological order.  The
resulting oriented broken path is denoted by \(C_\gamma\), and its terminal
point by
\[
    (A_\gamma,M_\gamma)
      =\left(\sum_{g_i=\mathsf A_{p_i}}p_i,
              \sum_{g_i=\mathsf M_{q_i}}q_i\right).
\]
\end{definition}

The terminal charge is insensitive to order, whereas the broken path generally
is not.  In particular, two histories may have the same ACS endpoint while
inducing different affine maps.  The missing information is recovered by the
following integral formula.

\begin{proposition}[Direct ACS evaluation]
\label{prop:acs-evaluation}
For every add--scale history \(\gamma\) and every initial value \(x\),
\begin{equation}
    \nu_x(\gamma)
      =e^{M_\gamma}
        \left(x+\int_{C_\gamma}e^{-M}\,dA\right).
    \label{eq:direct-acs-evaluation}
\end{equation}
\end{proposition}

\begin{proof}
Let an additive step \(\mathsf A_{p_i}\) occur when the cumulative
multiplicative charge is \(M_{i-1}\).  Along the corresponding horizontal edge
of \(C_\gamma\),
\[
    \int e^{-M}\,dA=p_i e^{-M_{i-1}}.
\]
After that addition, the remaining multiplicative steps have total charge
\(M_\gamma-M_{i-1}\).  Consequently the contribution of \(p_i\) to the final
value is
\[
    p_i e^{M_\gamma-M_{i-1}}
      =e^{M_\gamma}p_i e^{-M_{i-1}}.
\]
The initial value contributes \(e^{M_\gamma}x\), and vertical ACS edges
contribute nothing to the \(dA\)-integral.  Summing over the additive steps
gives \eqref{eq:direct-acs-evaluation}.  The argument uses oriented edge
integrals, so it also covers negative increments.
\end{proof}

The kernel \(e^{-M}\) is therefore forced by the affine cocycle: it normalizes
an additive contribution to the source frame, before the common terminal scale
\(e^{M_\gamma}\) returns it to the target frame.

\subsection{The history group and the descent of the charges}
\label{subsec:p0-history-group}

The preceding subsections attached charges to words.  To say precisely what the
ACS is, and in particular what it is *not*, it is useful to say which words are
being identified.  The identifications used below are only the ones the paper
has already made implicitly.

\begin{convention}[The history group]
\label{conv:p0-history-group}
Let
\[
  G=(\mathbb R,+)*(\mathbb R,+)
  =\bigl\langle \mathsf A_p,\ \mathsf M_q\ \bigm|\
    \mathsf A_p\mathsf A_{p'}=\mathsf A_{p+p'},\
    \mathsf M_q\mathsf M_{q'}=\mathsf M_{q+q'},\
    \mathsf A_0=\mathsf M_0=1 \bigr\rangle
\]
be the free product of two copies of the additive group of $\mathbb R$, written
multiplicatively, with generators $\mathsf A_p$ (additive steps) and
$\mathsf M_q$ (multiplicative steps).  A \emph{history} is an element of $G$
together with a chosen word representing it.  Two words differing only by
regrouping consecutive steps of one family represent the same element of $G$ and
have the same charge chain and the same evaluation, so nothing is lost by this
convention; words which interleave the two families are *not* identified.
\end{convention}

\begin{proposition}[Charge endpoint and evaluation are homomorphisms]
\label{prop:p0-history-homomorphisms}
The charge endpoint and the evaluation extend uniquely to homomorphisms
\[
  c:G\longrightarrow\mathbb R^2,
  \qquad
  \mathsf A_p\longmapsto(p,0),\quad \mathsf M_q\longmapsto(0,q),
\]
\[
  \nu:G\longrightarrow\operatorname{Aff}^{+}(1,\mathbb R),
  \qquad
  \mathsf A_p\longmapsto T_p,\quad \mathsf M_q\longmapsto D_{e^{q}},
\]
where $T_p(z)=z+p$ and $D_k(z)=kz$.  Both are surjective.  Moreover the charge
chain $C_\gamma$ is a function of the element of $G$ and not only of the word:
regrouping consecutive steps of one family replaces consecutive edges of one
charge path by the single edge with the same endpoints.
\end{proposition}

\begin{proof}
Each family is generated by elements satisfying the only relations imposed within
it, so a homomorphism is determined by its values on the generators and is
well defined; the images stated are themselves homomorphisms of $(\mathbb R,+)$.
Surjectivity of $c$ is clear, and surjectivity of $\nu$ holds because
translations and positive dilations generate
$\operatorname{Aff}^{+}(1,\mathbb R)$.  For the last claim, a regrouping
$\mathsf A_p\mathsf A_{p'}\mapsto\mathsf A_{p+p'}$ joins two consecutive
horizontal edges of $C_\gamma$ into the horizontal edge with the same endpoints
and the same orientation, and the vertical case is identical.
\end{proof}

\begin{proposition}[The ACS endpoint is the abelianization of the history group]
\label{prop:p0-acs-is-abelianization}
The charge homomorphism $c$ is the abelianization of $G$: its kernel is the
commutator subgroup,
\[
  \ker c=[G,G],
  \qquad\text{so}\qquad
  G_{\mathrm{ab}}\;\cong\;\mathbb R^2
\]
and the charge plane of Definition~\ref{def:acs} is the abelianization of the
history group.
\end{proposition}

\begin{proof}
For any groups $F,H$ one has
$(F*H)_{\mathrm{ab}}\cong F_{\mathrm{ab}}\oplus H_{\mathrm{ab}}$, because
abelianization is left adjoint to the inclusion of abelian groups and therefore
preserves coproducts.  With $F=H=(\mathbb R,+)$, which is abelian, this gives
$G_{\mathrm{ab}}\cong\mathbb R\oplus\mathbb R$.  Since $c$ is a surjection onto
the abelian group $\mathbb R^2$, it factors through $G_{\mathrm{ab}}$ as an
isomorphism, so $\ker c=[G,G]$.
\end{proof}

Now the question the review of this paper raised: does the charge pass to the
*evaluated* motion?  The next theorem answers it, and the answer is asymmetric.

\begin{theorem}[Only the multiplicative charge descends]
\label{thm:p0-descent}
Fix $q\ne0$ and $p\ne0$ and put
\begin{equation}\label{eq:p0-transport-element}
  h(p,q)=\mathsf M_q\,\mathsf A_p\,\mathsf M_{-q}\,\mathsf A_{-e^{-q}p}
  \;\in\; G .
\end{equation}
Then:
\begin{enumerate}
  \item $h(p,q)$ is a nontrivial element of $\ker\nu$; that is,
  $h(p,q)$ evaluates to the identity motion while being a nonempty history;
  \item its charges are
  \[
    c\bigl(h(p,q)\bigr)=\bigl(p(1-e^{-q}),\,0\bigr);
  \]
  \item the multiplicative charge $\mathsf M:G\to\mathbb R$,
  $\mathsf M=\mathrm{pr}_2\circ c$, factors through $\nu$: it equals
  $\log\circ L\circ\nu$, where $L:\operatorname{Aff}^{+}(1,\mathbb R)\to
  \mathbb R_{>0}$ is the linear part $(z\mapsto \Phi z+\xi)\mapsto\Phi$;
  \item the additive charge $\mathsf A:G\to\mathbb R$,
  $\mathsf A=\mathrm{pr}_1\circ c$, does *not* factor through $\nu$: there is no
  homomorphism $\varphi:\operatorname{Aff}^{+}(1,\mathbb R)\to\mathbb R^2$ with
  $c=\varphi\circ\nu$.  Equivalently, the additive charge is a property of the
  history and not of the motion it induces.
\end{enumerate}
\end{theorem}

\begin{proof}
(1) Apply the motions of $h(p,q)$ in chronological order to an arbitrary $z$:
\[
  z\ \xmapsto{\ \mathsf M_q\ }\ e^{q}z
    \ \xmapsto{\ \mathsf A_p\ }\ e^{q}z+p
    \ \xmapsto{\ \mathsf M_{-q}\ }\ z+e^{-q}p
    \ \xmapsto{\ \mathsf A_{-e^{-q}p}\ }\ z .
\]
Hence $\nu(h(p,q))=\mathrm{id}$.  It is nontrivial in $G$ because a reduced word
of a free product is trivial only if it is empty, and $q\ne0$, $p\ne0$ make
\eqref{eq:p0-transport-element} reduced.

(2) Add the charge increments in the same chronological order:
\[
  (0,q)+(p,0)+(0,-q)+(-e^{-q}p,0)=\bigl(p(1-e^{-q}),\,0\bigr).
\]

(3) $L$ is a homomorphism, and on generators $\log L(T_p)=0=\mathsf M(\mathsf A_p)$
and $\log L(D_{e^{q}})=q=\mathsf M(\mathsf M_q)$; since both sides are
homomorphisms on $G$ agreeing on generators, they agree on $G$.

(4) If $c=\varphi\circ\nu$ for some homomorphism $\varphi$ into the abelian group
$\mathbb R^2$, then $c$ would vanish on $\ker\nu$.  But by (1) and (2),
$h(p,q)\in\ker\nu$ while $c(h(p,q))=(p(1-e^{-q}),0)\ne(0,0)$ for $p\ne0$,
$q\ne0$.  Contradiction.  (The structural reason is
Proposition~\ref{prop:p0-acs-is-abelianization} together with the computation
$\operatorname{Aff}^{+}(1,\mathbb R)_{\mathrm{ab}}\cong\mathbb R_{>0}$ recorded
in the opening of this section: every homomorphism from the motion group to an
abelian group is one-dimensional, whereas the charges are two-dimensional.  The
explicit element $h(p,q)$ is what makes the failure a proof rather than a
dimension count.)
\end{proof}

Three consequences deserve to be stated together, because they are the precise
version of statements that earlier drafts left informal.

\begin{enumerate}
  \item \emph{The ACS is the commutativization of the history language, and the
  sense of that phrase is now fixed.}  It is the abelianization of $G$
  (Proposition~\ref{prop:p0-acs-is-abelianization}), not of the motion group.
  The motion group's abelianization is one-dimensional and keeps the scale alone;
  the additive charge does not survive evaluation at all
  (Theorem~\ref{thm:p0-descent}).
  \item \emph{What the ACS retains beyond its abelianized endpoint is the path.}
  Two histories with the same charge endpoint may have different charge chains;
  the elementary rectangle $\gamma=(\mathsf A_p,\mathsf M_q)$ and
  $\delta=(\mathsf M_q,\mathsf A_p)$ is the smallest example, with equal
  endpoints $(p,q)$ and different chains.  Proposition~\ref{prop:acs-evaluation}
  uses the chain, not the endpoint, and that is why it recovers the evaluated
  motion from commutative data.  A charge *endpoint* observable therefore
  cannot see the difference, exactly as
  Proposition~\ref{prop:p0-characters-cannot-see-tearing} states, while the
  endpoint together with the path can.
  \item \emph{The two failure directions are dual.}  Histories with equal motion
  and different charges are the elements of $\ker\nu$ of
  Theorem~\ref{thm:p0-descent}; histories with equal charges and different
  motions are the pairs measured by the relative torsion of
  Definition~\ref{def:relative-torsion}.  The first failure is why the charge map
  cannot be pushed through evaluation; the second is why the evaluation cannot be
  recovered from the charge endpoint.  The ACS is the object in which both become
  computable, because it carries the endpoint (first failure visible) and the
  path (second failure measured by the weighted area of
  Theorem~\ref{thm:torsion-stokes}).
\end{enumerate}

\subsection{Compatibility and relative torsion}

There are two useful compatibility conditions, and they should not be
conflated.  Histories \(\gamma\) and \(\delta\) are \emph{scale-compatible} if
\[
    M_\gamma=M_\delta.
\]
They are \emph{charge-compatible} if, in addition,
\[
    A_\gamma=A_\delta.
\]
Scale compatibility makes their affine maps have the same linear part.  Charge
compatibility makes their ACS paths have both endpoints in common and hence
makes \(C_\gamma-C_\delta\) a closed oriented one-chain.

\begin{definition}[Relative arithmetic torsion]
\label{def:relative-torsion}
For scale-compatible add--scale histories \(\gamma\) and \(\delta\), their \emph{relative
arithmetic torsion} is
\begin{equation}
    \tau(\gamma,\delta)
       :=\xi_\gamma-\xi_\delta,
    \label{eq:relative-torsion}
\end{equation}
where \(\xi_\gamma\) and \(\xi_\delta\) are the target-frame translation
terms of the two affine operators.  This relative translation defect is the
primary torsion invariant used below.
\end{definition}

\begin{proposition}
\label{prop:relative-torsion-independent}
For every initial value \(x\), relative arithmetic torsion equals the endpoint
difference
\[
    \tau(\gamma,\delta)
      =\nu_x(\gamma)-\nu_x(\delta),
\]
and is therefore independent of \(x\).
If \(M_\gamma=M_\delta=M_*\), then
\begin{equation}
    \tau(\gamma,\delta)
      =e^{M_*}
        \left(
          \int_{C_\gamma}e^{-M}\,dA
          -\int_{C_\delta}e^{-M}\,dA
        \right).
    \label{eq:relative-torsion-boundary-preparation}
\end{equation}
\end{proposition}

\begin{proof}
Apply Proposition~\ref{prop:acs-evaluation} to both histories.  Their initial
terms are respectively \(e^{M_\gamma}x\) and \(e^{M_\delta}x\), which cancel
under scale compatibility.  The remaining terms give
\eqref{eq:relative-torsion-boundary-preparation}.
\end{proof}

Temporal reversal gives one possible comparison, but it is not built into the
definition.  The pair \((\gamma,\delta)\) may consist of any two
scale-compatible add--scale histories.  This extension is important because order defects
occur between histories that are neither reversals nor planar mirrors of one
another.

\subsection{The torsion--Stokes formula}

Suppose now that \(\gamma\) and \(\delta\) are charge-compatible, with common
terminal charge \((A_*,M_*)\).  Define the target-normalized one-form
\begin{equation}
    \eta_*:=e^{M_*-M}\,dA
    \label{eq:target-normalized-acs-form}
\end{equation}
on the ACS.  Relative torsion is then the integral of \(\eta_*\) over the
closed one-chain \(C_\gamma-C_\delta\).  The chosen order of the two paths
fixes the sign.

\begin{theorem}[Weighted torsion--Stokes theorem]
\label{thm:torsion-stokes}
Let \(\gamma\) and \(\delta\) be charge-compatible add--scale histories with
common terminal multiplicative charge \(M_*\).  Let
\(\Sigma_{\gamma,\delta}\) be any oriented piecewise smooth singular
two-chain in the ACS satisfying
\[
    \partial\Sigma_{\gamma,\delta}=C_\gamma-C_\delta.
\]
Then
\begin{align}
    \tau(\gamma,\delta)
      &=\int_{C_\gamma-C_\delta}\eta_*
        \label{eq:torsion-boundary-integral}\\
      &=\int_{\Sigma_{\gamma,\delta}}
          e^{M_*-M}\,dA\wedge dM.
        \label{eq:torsion-area-integral}
\end{align}
The value of the area integral is independent of the chosen filling chain.
\end{theorem}

\begin{proof}
Equation \eqref{eq:relative-torsion-boundary-preparation} and charge
compatibility give
\[
    \tau(\gamma,\delta)
       =\int_{C_\gamma-C_\delta}e^{M_*-M}\,dA
       =\int_{C_\gamma-C_\delta}\eta_*.
\]
With the orientation \(dA\wedge dM\), direct differentiation yields
\[
    d\eta_*
      =d(e^{M_*-M})\wedge dA
      =-e^{M_*-M}dM\wedge dA
      =e^{M_*-M}dA\wedge dM.
\]
Stokes' theorem for piecewise smooth chains therefore gives
\[
    \int_{C_\gamma-C_\delta}\eta_*
      =\int_{\partial\Sigma_{\gamma,\delta}}\eta_*
      =\int_{\Sigma_{\gamma,\delta}}d\eta_*,
\]
which is \eqref{eq:torsion-area-integral}.  Such a filling exists because the
ACS plane is contractible.  Finally, if \(\Sigma\) and \(\Sigma'\) have the
same boundary, applying Stokes' theorem to each shows
\[
    \int_\Sigma d\eta_*
       =\int_{\partial\Sigma}\eta_*
       =\int_{\partial\Sigma'}\eta_*
       =\int_{\Sigma'}d\eta_*.
\]
This also proves filling independence.  The chain formulation remains valid
when the broken paths self-intersect or carry signed increments; no informal
notion of an enclosed region is required.
\end{proof}

Charge compatibility is needed for the closed-chain and area formulas, but not
for Definition~\ref{def:relative-torsion}.  A merely scale-compatible pair
still has a well-defined translation defect, even when the endpoints of its ACS
paths have different \(A\)-coordinates.

Figure~\ref{fig:acs-compatible-paths-weighted-filling} depicts the chain
orientation used in Theorem~\ref{thm:torsion-stokes}.  It is a schematic of
the theorem's data rather than an argument for filling independence.

\begin{figure}[t]
\centering
\begin{tikzpicture}[
  x=1.15cm,y=1.05cm,
  axis/.style={-{Latex[length=2mm]},thick,black!70},
  gamma/.style={-{Latex[length=2mm]},very thick,red!65!black},
  delta/.style={-{Latex[length=2mm]},very thick,blue!65!black},
  every node/.style={font=\small}
]
  \coordinate (O) at (0,0);
  \coordinate (g1) at (2.2,0);
  \coordinate (g2) at (2.2,0.9);
  \coordinate (g3) at (4,0.9);
  \coordinate (P) at (4,2.7);
  \coordinate (d1) at (0,1.5);
  \coordinate (d2) at (1.3,1.5);
  \coordinate (d3) at (1.3,2.7);

  \path[fill=violet!11,draw=none]
    (O)--(g1)--(g2)--(g3)--(P)--(d3)--(d2)--(d1)--cycle;
  \draw[axis] (-0.25,0) -- (4.75,0) node[right] {$A$};
  \draw[axis] (0,-0.25) -- (0,3.25) node[above] {$M$};

  \draw[gamma] (O)--(g1);
  \draw[gamma] (g1)--(g2);
  \draw[gamma] (g2)--(g3);
  \draw[gamma] (g3)--(P);
  \draw[delta] (O)--(d1);
  \draw[delta] (d1)--(d2);
  \draw[delta] (d2)--(d3);
  \draw[delta] (d3)--(P);

  \fill (O) circle (1.7pt) node[below left] {$(0,0)$};
  \fill (P) circle (1.7pt) node[above right] {$(A_*,M_*)$};
  \node[red!65!black,fill=white,inner sep=1pt] at (3.3,0.72) {$C_\gamma$};
  \node[blue!65!black,fill=white,inner sep=1pt] at (0.72,1.75) {$C_\delta$};
  \node[align=center,violet!55!black] at (2.55,1.8)
    {$\Sigma_{\gamma,\delta}$\\
     $e^{M_*-M}\,dA\wedge dM$};
  \node[font=\scriptsize,anchor=west] at (4.35,1.72)
    {$\partial\Sigma=C_\gamma-C_\delta$};

  \draw[-{Latex[length=1.7mm]},black!55] (4.55,0.35)--(4.95,0.35)
    node[right,font=\scriptsize] {$dA$};
  \draw[-{Latex[length=1.7mm]},black!55] (4.55,0.35)--(4.55,0.8)
    node[above,font=\scriptsize] {$dM$};
\end{tikzpicture}
\caption{Two charge-compatible ACS histories.  Both paths run from $(0,0)$
to the common charge endpoint $(A_*,M_*)$.  Traversing $C_\gamma$ forward and
$C_\delta$ backward gives the positive boundary of the displayed filling for
the orientation $dA\wedge dM$; the density is weighted by $e^{M_*-M}$.}
\label{fig:acs-compatible-paths-weighted-filling}
\end{figure}

\subsection{Explicit order comparisons}

\paragraph{The elementary rectangle.}
Let
\[
    \gamma=(\mathsf A_p,\mathsf M_q),
    \qquad
    \delta=(\mathsf M_q,\mathsf A_p).
\]
The path \(C_\gamma-C_\delta\) is the positively oriented boundary of the
rectangle \([0,p]\times[0,q]\) when \(p,q>0\).  Direct evaluation and the
weighted-area formula agree:
\begin{align*}
    \tau(\gamma,\delta)
      &=e^q(x+p)-(e^q x+p)\\
      &=p(e^q-1)\\
      &=\int_0^p\int_0^q e^{q-M}\,dM\,dA.
\end{align*}
Thus ``add, then multiply'' minus ``multiply, then add'' is positive under the
orientation fixed in Definition~\ref{def:acs}.  In the flow notation used in
the next chapter, \(p=\mu h\) and \(q=\lambda k\), so the defect is
\[
    \mu h(e^{\lambda k}-1).
\]

\paragraph{A four-step staircase.}
For
\[
    \gamma=(\mathsf A_p,\mathsf M_r,\mathsf A_q,\mathsf M_s),
    \qquad
    \delta=(\mathsf M_r,\mathsf A_p,\mathsf M_s,\mathsf A_q),
\]
the histories have the same total charge \((p+q,r+s)\), but \(\delta\) is not
the temporal reversal of \(\gamma\).  Their evaluations are
\[
    \nu_x(\gamma)=e^{r+s}x+p e^{r+s}+q e^s,
\]
\[
    \nu_x(\delta)=e^{r+s}x+p e^s+q.
\]
Consequently
\begin{equation}
    \tau(\gamma,\delta)
      =p e^s(e^r-1)+q(e^s-1).
    \label{eq:four-step-staircase-torsion}
\end{equation}
For positive \(p,q,r,s\), both terms are positive weighted strips in an ACS
filling.  Formula \eqref{eq:four-step-staircase-torsion} also shows how a later
scale changes the weight of an earlier commutation defect.

\paragraph{Equal charges without reversal.}
The pair
\[
    (\mathsf A_p,\mathsf M_r,\mathsf A_q),
    \qquad
    (\mathsf A_q,\mathsf M_r,\mathsf A_p)
\]
has common charge \((p+q,r)\), and its relative torsion is
\[
    (p-q)(e^r-1).
\]
It may be positive, negative, or zero.  In particular, equal charge does not
determine torsion, while zero torsion does not imply equality of marked
histories.

\subsection{What the ACS forgets}
\label{subsec:p0-acs-forgets}

At the formal level, the map
\[
    \gamma\longmapsto(A_\gamma,M_\gamma)
\]
is an order-forgetting charge homomorphism from concatenated marked histories
to the additive group \(\mathbb R^2\).  It is therefore analogous to an
abelianization map.  This statement concerns the formal charge language only.
It does not identify the symbolic history space with the evaluated affine
group, and it does not assert that either object is a free group without an
additional choice of formal generators and relations.

The content of Theorem~\ref{thm:torsion-stokes} is precisely what survives this
loss of order: the endpoint charges specify a common frame, while the weighted
filling measures the translation defect between two histories.
Section~\ref{sec:p0-contact} retains the current value as a third coordinate and
identifies the infinitesimal form of the same defect with contact curvature.

\subsection{Tearing}
\label{subsec:p0-tearing}

We can now name the phenomenon that the three preceding chapters have been
describing from three sides.  The name is new here, the mathematics is not: it
is the content of Definition~\ref{def:relative-torsion} and
Theorem~\ref{thm:torsion-stokes} read geometrically.

\begin{definition}[Tearing of arithmetic motion]
\label{def:p0-tearing}
Let $\gamma$ and $\delta$ be two add--scale histories of a regular AES that are
scale-compatible in the sense of Section~\ref{sec:p0-acs-tearing}, and let
$\tau(\gamma,\delta)$ be their relative torsion.  The pair is said to \emph{tear}
the AES by $\tau(\gamma,\delta)$: the two histories induce the same linear part,
hence agree on the multiplicative content of the motion, while their additive
content differs by exactly $\tau(\gamma,\delta)$.
\end{definition}

The identification of that scalar defect with an area requires a second
hypothesis, and the two must not be merged.  By
Theorem~\ref{thm:torsion-stokes}, if $\gamma$ and $\delta$ are in addition
*charge-compatible*, then $\tau(\gamma,\delta)$ equals the integral of the
weighted area form over any oriented filling of the closed chain
$C_\gamma-C_\delta$.  For a merely scale-compatible pair the chain
$C_\gamma-C_\delta$ is not closed, because the two paths have different
additive endpoints, so there is no such filling; the scalar defect
\eqref{eq:relative-torsion} remains defined, but it is not yet an area.

Three properties justify treating tearing as a geometric notion rather than a
bookkeeping defect.

\begin{enumerate}
  \item \emph{It exists only relative to the commutativization.}  In the AES
  itself, $\gamma$ and $\delta$ are simply two motions; there is no region
  bounded by them, because the AES is where the two rulers fail to commute.  The
  region appears only after passing to the ACS, and the weight of the filling is
  forced by the evaluation formula \eqref{eq:direct-acs-evaluation}.
  \item \emph{For charge-compatible pairs it is measured, not merely detected.}
  There $\tau(\gamma,\delta)$ is an exact weighted area
  \eqref{eq:torsion-area-integral}; for the elementary rectangle it is
  $\mu h(e^{\lambda k}-1)$, which is the discrete Baumslag--Solitar transport of
  \eqref{eq:p0-bs-relation} read as a defect.
  \item \emph{It is infinitesimally rigid.}  Its density is the constant
  $\mu\lambda$ of \eqref{eq:infinitesimal-order-defect}, which is also the
  bracket coefficient of the arithmetic frame
  \eqref{eq:arithmetic-frame-bracket-on-a}.  Section~\ref{sec:p0-contact}
  shows that this coefficient is the curvature of the specified horizontal
  connection, hence
  that tearing cannot be removed by any local recalibration of the two
  arithmetic directions in which both remain horizontal.
\end{enumerate}

The word \emph{tearing} is chosen because the two-dimensional picture is that of
a surface cut along the two motions: an additive edge and a multiplicative edge
that close up on the level of total charge but not on the level of motion.  The
standard description of a rip in a material is also a pair of paths that agree
far away and separate inside; the arithmetic version of that picture is
Definition~\ref{def:p0-tearing}.  As with ripple geometry, the terminology is
descriptive and asserts nothing beyond the displayed identities; what it adds is
a single name for the phenomenon whose exact measure is the ACS area, whose
infinitesimal density is $\mu\lambda$, and whose carrier in a punctured model is
the subject of Section~\ref{sec:p0-holed}.
